% ============================================================
%  CHUONG 5. KET LUAN VA HUONG PHAT TRIEN
% ============================================================
\chapter{Kết luận và hướng phát triển}
\label{ch:ket-luan}

\section{Kết luận}

Đồ án đã hoàn thành việc xây dựng \textbf{CookingAdvisor}, một ứng dụng web hỗ
trợ người nấu ăn trong gia đình ở ba khâu: tìm món từ nguyên liệu sẵn có, lập
thực đơn cho cả tuần và tổng hợp danh sách đi chợ.

\subsection{Kết quả đạt được}

\textbf{Về mặt chức năng,} hệ thống cung cấp 20 chức năng phân thành ba nhóm theo
mức quyền truy cập, chạy hoàn chỉnh đầu-cuối trên cơ sở dữ liệu mẫu gồm 26 món ăn
Việt Nam, 40 nguyên liệu và 10 danh mục. Toàn bộ sáu mục tiêu đặt ra ở phần Mở
đầu đều đã đạt, với kết quả kiểm chứng cụ thể trình bày ở
mục~\ref{sec:ket-qua-dat-duoc}.

\textbf{Về mặt thuật toán,} ba thuật toán lõi đã được thiết kế, cài đặt và kiểm
chứng:

\begin{itemize}
  \item \textit{Thuật toán gợi ý theo nguyên liệu} sử dụng độ phủ tập hợp kết hợp
        xếp hạng bốn tiêu chí. Đóng góp đáng chú ý ở đây là lập luận chọn độ phủ
        thay cho hệ số Jaccard: do Jaccard là độ đo đối xứng nên nó phạt cả phần
        nguyên liệu dư của người dùng, dẫn tới việc một món nấu được ngay có thể
        bị xếp hạng rất thấp.
  \item \textit{Thuật toán sinh thực đơn tuần} theo hướng tham lam có ràng buộc,
        với thứ tự nới lỏng ràng buộc được quy định rõ ràng và ràng buộc vùng
        miền được giữ cứng.
  \item \textit{Thuật toán sinh danh sách đi chợ} gộp nhóm theo khóa ghép (nguyên
        liệu, đơn vị), bảo đảm không cộng dồn nhầm hai giá trị khác đơn vị.
\end{itemize}

\textbf{Về mặt kiểm chứng,} bộ 29 ca kiểm thử đơn vị đã phát hiện được hai khiếm
khuyết mà việc quan sát bằng mắt thường khó nhận ra:

\begin{enumerate}[label=\arabic*.]
  \item \textit{Suy thoái tính đa dạng của thực đơn.} Khi thư viện món không đủ
        lớn, tiêu chí phụ theo năng lượng khiến thuật toán chọn lại đúng một món
        cho hầu hết các suất còn trống, đo được là một món chiếm 12 suất. Sau khi
        bổ sung tiêu chí rải đều theo số lần đã dùng, thực đơn đạt 17 món khác
        nhau trên 21 suất.
  \item \textit{Thứ tự không xác định khi phân trang.} Việc sắp xếp chỉ theo một
        khóa duy nhất khiến các bản ghi có giá trị bằng nhau không có thứ tự ổn
        định, dẫn tới nguy cơ một món xuất hiện ở hai trang. Đã khắc phục bằng
        khóa sắp xếp phụ theo tên món.
\end{enumerate}

Việc phát hiện được hai khiếm khuyết này là minh chứng cho giá trị thực tế của
kiểm thử tự động, đồng thời cho thấy quyết định kiến trúc tách tầng Service khỏi
tầng Controller (mục~\ref{sec:kien-truc}) đã phát huy tác dụng: nếu logic nằm
trong Controller, việc viết các ca kiểm thử này sẽ khó khăn hơn nhiều.

\textbf{Về mặt giao diện,} hệ thống hiển thị đúng trên dải màn hình từ 375 px tới
1440 px, không xuất hiện cuộn ngang, hỗ trợ chế độ tối và đạt chuẩn tương phản
WCAG 2.1 mức AA với cặp màu thấp nhất đo được là 4,69:1.

\subsection{Đóng góp của đồ án}

So với các ứng dụng đã khảo sát ở mục~\ref{sec:khao-sat}, đồ án có ba điểm khác
biệt.

\textbf{Thứ nhất, tích hợp trọn vẹn ba chức năng trong một luồng công việc liền
mạch.} Các sản phẩm khảo sát thường mạnh ở một khâu và thiếu ở khâu khác. Mealime
lập được thực đơn nhưng không hỗ trợ truy vấn ngược từ nguyên liệu. Paprika quản
lý công thức và danh sách đi chợ tốt nhưng không tự động hóa việc lập thực đơn.
Các trang trong nước chủ yếu dừng ở việc cung cấp nội dung.

\textbf{Thứ hai, dữ liệu món ăn Việt Nam được chuẩn hóa về định lượng.} Bảng
trung gian giữa món ăn và nguyên liệu mang thêm thuộc tính khối lượng và đơn vị.
Đây là điều kiện tiên quyết để danh sách đi chợ có thể cộng dồn chính xác, điều
mà các nền tảng có nội dung do cộng đồng đóng góp khó bảo đảm.

\textbf{Thứ ba, thuật toán minh bạch và kiểm chứng được.} Toàn bộ tiêu chí xếp
hạng và thứ tự nới lỏng ràng buộc đều được phát biểu tường minh, có mã giả, có
lưu đồ và có ca kiểm thử tương ứng.

\subsection{Hạn chế}

Đồ án còn tám hạn chế đã nêu chi tiết ở mục~\ref{sec:han-che}. Trong đó, ba hạn
chế có ảnh hưởng lớn nhất tới trải nghiệm người dùng là:

\begin{itemize}
  \item Thuật toán gợi ý mới xét việc \textbf{có hay không có} một nguyên liệu,
        chưa xét người dùng có \textbf{đủ khối lượng} hay không.
  \item Kết quả gợi ý \textbf{giống nhau với mọi người dùng}, chưa khai thác lịch
        sử nấu ăn để cá nhân hóa.
  \item Thông tin dinh dưỡng mới dừng ở \textbf{năng lượng theo khẩu phần}, chưa
        bóc tách các thành phần dinh dưỡng chi tiết.
\end{itemize}

\section{Hướng phát triển}

\subsection{Hoàn thiện các thuật toán hiện có}

\textbf{Gợi ý có xét định lượng.} Mở rộng mô hình dữ liệu để người dùng khai báo
được khối lượng nguyên liệu đang có, không chỉ khai báo có hay không. Khi đó điều
kiện ``nấu được ngay'' trở thành: với mọi nguyên liệu $i$ của món, khối lượng
người dùng có phải lớn hơn hoặc bằng khối lượng công thức yêu cầu. Thay đổi này
cũng cho phép danh sách đi chợ \textbf{trừ đi phần đã có}, chỉ liệt kê phần còn
thiếu.

\textbf{Lập thực đơn theo hướng tối ưu.} Thay chiến lược tham lam bằng một thuật
toán tìm kiếm cục bộ hoặc quy hoạch ràng buộc, để tổng năng lượng theo ngày bám
sát mục tiêu hơn. Có thể bổ sung các ràng buộc mới như cân đối nhóm thực phẩm
trong tuần hoặc tránh lặp cùng một nguyên liệu chính trong hai ngày liên tiếp.

\textbf{Cá nhân hóa kết quả gợi ý.} Khi hệ thống đã tích lũy đủ dữ liệu về lịch
sử nấu ăn và món yêu thích, có thể áp dụng kỹ thuật lọc cộng tác để bổ sung một
tiêu chí xếp hạng dựa trên mức độ phù hợp với khẩu vị của từng người dùng. Cần
lưu ý đây là hướng chỉ khả thi khi đã có người dùng thực, do các phương pháp này
gặp vấn đề khởi đầu nguội trên hệ thống mới.

\subsection{Mở rộng chức năng}

\begin{table}[H]
\centering
\caption{Các hướng mở rộng chức năng}
\label{tab:mo-rong}
\begin{tabularx}{\textwidth}{p{4.2cm} Y}
\toprule
\bfseries Hướng mở rộng & \bfseries Nội dung \\
\midrule
Dinh dưỡng chi tiết & Bổ sung đạm, béo, tinh bột, chất xơ; cảnh báo theo chế độ
  ăn hoặc bệnh lý \\
\rowhl Quản lý tài khoản & Đổi mật khẩu, quên mật khẩu, xác thực email, trang hồ
  sơ cá nhân \\
Tải ảnh lên & Thay việc nhập đường dẫn ảnh bằng chức năng tải tệp và tự tạo ảnh
  thu nhỏ \\
\rowhl Cộng đồng & Cho phép người dùng đóng góp công thức, đánh giá và bình
  luận \\
Chia sẻ thực đơn & Xuất thực đơn tuần ra tệp hoặc chia sẻ qua đường dẫn công
  khai \\
\rowhl In danh sách đi chợ & Xuất bản in hoặc tệp PDF để mang theo khi đi chợ \\
\bottomrule
\end{tabularx}
\end{table}

\subsection{Cải thiện kỹ thuật}

\textbf{Phân trang và tìm kiếm cho khu quản trị.} Hiện các danh sách trong khu
quản trị hiển thị toàn bộ bản ghi. Với quy mô dữ liệu mẫu thì chưa gây trở ngại,
nhưng sẽ trở thành hạn chế khi số món tăng lên.

\textbf{Bổ sung kiểm thử tích hợp và kiểm thử giao diện tự động.} Hiện phần giao
diện mới được kiểm tra thủ công. Có thể bổ sung kiểm thử tích hợp cho các luồng
Controller và kiểm thử giao diện tự động cho các kịch bản chính. Việc này cũng
giúp phủ được nhóm hành vi mà nhà cung cấp EF Core In-Memory không mô phỏng được,
như đã nêu ở mục~\ref{sec:kiem-thu}.

\textbf{Ứng dụng di động.} Bối cảnh sử dụng thực tế của chức năng danh sách đi
chợ là khi người dùng đang ở chợ hoặc siêu thị. Một ứng dụng di động có khả năng
hoạt động ngoại tuyến sẽ phù hợp hơn giao diện web.
